% ---------------------------------------------------------------------------
% Electronic supplementary material for
% "Payoff scale reshapes how language models play the prisoner's dilemma".
%
%   pdflatex supplementary && bibtex supplementary && pdflatex supplementary
%
% All tables are machine-generated into supp_tables_auto.tex by
% Analysis/scaling/s07_supplementary.py.  Do not retype numbers here.
% ---------------------------------------------------------------------------
\documentclass[11pt,a4paper]{article}

\usepackage[T1]{fontenc}
\usepackage[utf8]{inputenc}
\usepackage{lmodern}
\usepackage[margin=1in]{geometry}
\usepackage{graphicx}
\usepackage{amsmath,amssymb}
\usepackage{booktabs}
\usepackage{caption}
\usepackage[numbers,sort&compress]{natbib}
\usepackage{microtype}
\usepackage{url}\urlstyle{tt}   % \path breaks long file paths at the slashes
\usepackage{xcolor}
\usepackage[colorlinks=true,allcolors=blue!55!black]{hyperref}
\setlength{\emergencystretch}{2em}

\captionsetup{font=small,labelfont=bf,width=\textwidth}
\renewcommand{\baselinestretch}{1.12}
\renewcommand{\thetable}{S\arabic{table}}
\renewcommand{\thefigure}{S\arabic{figure}}
\renewcommand{\thesection}{S\arabic{section}}

\newcommand{\lam}{\lambda}

\date{}

\begin{document}

\begin{center}
{\LARGE\bfseries Electronic supplementary material}\\[6pt]
{\large Payoff scale reshapes how language models play\\[2pt]
the prisoner's dilemma}
\end{center}
\bigskip

\tableofcontents
\bigskip

\section{Scope of the corpus}

The analysis in the main text uses one corpus of language-model transcripts,
\path{Dataset/data_fairgame_frontier_llm}, and one synthetic corpus,
\path{Dataset/noise_dataset}, which trains the LSTM branch of the strategy
read-out and contains no language-model output at all. No other dataset in the
repository enters any number reported in the paper.

The transcript corpus is a complete crossing of ten payoff scales, five models,
five prompt languages, four persona pairings and ten replicates, at ten rounds
per game and two agents per dyad. Table~\ref{tab:S-corpus} gives its
composition. Every one of the 250 model-by-language-by-scale cells holds
exactly 40 dyads per language, or 200 per cell of the coarser
model-by-scale grid, and the ingest raises an error rather than writing its
tables if that is not what is on disk. It also checks that every game is ten
rounds long, and that the payoff scale recovered from the recorded payoffs
agrees with the scale the cell was run at, which catches a mislabelled
directory.

\section{Cooperation, by scale, model and language}

Table~\ref{tab:S-scale} gives the cooperation rate behind
Fig.~1b of the main text, with bootstrap intervals.
Table~\ref{tab:S-language} breaks the same quantity down by prompt language and
is the table behind Fig.~2.

The final column of Table~\ref{tab:S-language} is the quantity we draw
attention to in the main text: the range of cooperation across the ten payoff
scales, computed within one model and one language. It runs from 0.080
(Gemini 3.1 Flash-Lite in Arabic) to 0.449 (GPT-5.4 Nano in English). The
sensitivity of a model to a transformation that changes nothing about the game
is therefore not a single number even for a fixed model; it depends on the
language the game was posed in.

\section{Variance decomposition}

Table~\ref{tab:S-variance} is the full analysis-of-variance table summarised in
the main text. It is computed on the 250 cell means rather than on the 20{,}000
agent-games, because the design is exactly balanced at the cell level, which
makes the type-II decomposition order-independent.

The reading that matters for benchmark practice is the contrast between rows.
The payoff scale carries $0.7\%$ of the variance across cells as a main effect
and is not significant ($p = 0.36$). Its interaction with the model carries
$6.5\%$ and its interaction with the prompt language $3.8\%$, and both are
significant. A design that estimates a main effect for the payoff scale,
pooling over models and languages, is therefore built to return exactly the
null result it would return if no model responded to the scale at all.

\section{Persona}

Table~\ref{tab:S-persona} gives the persona effect at every scale and model,
which is the table behind Fig.~3a. The two leftmost columns are the scales at
which every payoff printed is at most 1, since with base payoffs $0, 2, 6, 10$
the largest number shown is $10\lam$.

Two features are worth stating explicitly because they are easy to misread from
the figure. First, the sign reversal is a reversal of the \emph{persona} effect,
not of the cooperation rate: below the boundary the three reversing models
cooperate more when told they are selfish than when told they are cooperative.
Second, Qwen3 235B-A22B is inverted at every scale rather than only below the
boundary, cooperating 0.072 under the cooperative persona against 0.798 under
the selfish one. We report it because excluding it would flatter the analysis,
but the paper draws no conclusion from the inversion itself beyond noting that
its magnitude, like everything else measured here, depends on the payoff scale.

\section{Strategy read-out}

\subsection{What the two branches do}

Each agent-game is encoded as a sequence of two-letter tokens, one per round,
pairing the previous round's joint outcome with the action taken in the current
round. The alphabet is shared between the synthetic training corpus and the
language-model transcripts, which is what makes a classifier trained on the
former applicable to the latter without any adaptation.

The rule base asks an exact question: is this trajectory precisely what AllC,
AllD, Tit-for-Tat or Win-Stay-Lose-Shift would have played, given the history
that actually occurred? Because all four rules are memory-one, the answer is a
per-token lookup and requires no learning. It returns a set rather than a single
label, because over ten rounds several canonical rules routinely prescribe the
same moves. A player cooperating in all ten rounds against a cooperator is
exactly consistent with three of the four rules at once, AllC, Tit-for-Tat and
Win-Stay-Lose-Shift, and no method, learned or otherwise, can separate them on
that evidence. Reporting the set is the correct answer, not a failure.

The LSTM ranks within whatever set the rule base returns, and decides alone only
where the set is empty. It therefore cannot overturn a deduction; it can only
break a tie the evidence genuinely leaves open, or supply an attribution where
no canonical rule applies.

\subsection{Provenance, and which claims depend on the classifier}

Over the whole corpus, $23.0\%$ of agent-games are matched exactly by one
canonical rule, $13.2\%$ by several at once, and $63.7\%$ by none. Those three
provenances are reported separately throughout, and the claims in the main text
are ordered by how much they depend on the learned component:

\begin{itemize}
\item The share of agent-games matched exactly by one rule
      (Table~\ref{tab:S-deduced}) and the number of rounds violating the
      nearest canonical rule are computed by exact matching. Neither depends on
      the LSTM in any way. These carry the main text's strategy claim.
\item The compositional claim, that the mix moves from unconditional towards
      reciprocal play (Table~\ref{tab:S-mix}), uses the labels, and therefore
      leans on the LSTM for the $76.9\%$ of agent-games that are ambiguous or
      unmatched. It is the weakest of the three and is presented as such.
\end{itemize}

Table~\ref{tab:S-provenance} shows that this dependence is very unevenly spread
across the four labels, which bounds how the compositional result may be read.
AllC and AllD carry deduced shares of $26.0\%$ and $32.7\%$, so a movement in
those two is anchored in part by exact matching. Tit-for-Tat and
Win-Stay-Lose-Shift are $95.4\%$ and $97.3\%$ unmatched, meaning the label is an
attribution the LSTM makes about which canonical rule the play most resembles,
on trajectories that follow none of them exactly. A rise in the Tit-for-Tat
share is therefore evidence that play is drifting towards the reciprocal region
of the space, and is not evidence that any model is implementing Tit-for-Tat.
The main text states this where the result is reported.

\subsection{Validation}

Table~\ref{tab:S-classifier} reports the accuracy of the LSTM branch. Training
uses the noise-free and 5\% execution-noise variants of the synthetic corpus
together, which forces the network to identify a rule from imperfect behaviour
rather than from a clean signature; the 10\% and 20\% variants are held out
entirely and measure how the read-out degrades as play departs further from any
rule. Accuracy falls from $98.5\%$ in distribution to $93.6\%$ and $81.1\%$ at
those two unseen noise levels, which is the honest bound on how far the label
distributions in Table~\ref{tab:S-mix} should be pushed.

\section{Why the model curves cancel under pooling}

Table~\ref{tab:S-shapes} gives the pairwise correlations between the five
models' cooperation-against-scale curves. Six of the ten pairs are negative,
the mean is $-0.10$, and the extremes are $-0.72$ and $+0.73$. Averaging curves
that disagree in shape this strongly reduces the range from a mean of 0.136
within a model to 0.048 for the average, a factor of 2.8. The cancellation is
arithmetic rather than statistical, and no increase in sample size affects it.

\section{Reproduction}

The pipeline runs in order and each stage writes its outputs under
\path{Analysis/scaling/}:

\begin{itemize}
\item \path{s00_build.py} parses the transcripts into \texttt{rounds.parquet}
      and \texttt{games.parquet}, and refuses to write either if the grid is
      incomplete.
\item \path{s01_train_lstm.py} trains the LSTM branch on the synthetic
      corpus and writes \path{models/strategy_lstm.pt}.
\item \path{s02_readout.py} applies the hybrid read-out to all 20{,}000
      agent-games and writes \texttt{readout.parquet}.
\item \path{s03_stats.py} and \path{s04_strategy_stats.py} compute every
      statistic in the paper into \path{tables/}.
\item \path{s05_figures.py} draws Figs.~1 to 5.
\item \path{s06_tables.py} and \path{s07_supplementary.py} generate the
      LaTeX tables in the main text and in this supplement.
\end{itemize}

Random seeds are fixed: 1234 for classifier training, 20260905 for the
bootstrap and the permutation test.

\section{Supplementary tables}

% Generated by Analysis/scaling/s07_supplementary.py - do not edit.
\begin{table}[htbp]
\centering
\footnotesize
\setlength{\tabcolsep}{4pt}
\caption{\textbf{The corpus, by model.} Every model was run on the same ten payoff scales and five prompt languages. ``Cell $n$'' is the number of dyads in each model $\times$ language $\times$ scale cell; the design is balanced with none missing, and the build refuses to write its tables otherwise.}
\label{tab:S-corpus}
\begin{tabular}{@{}llccrrrc@{}}
\toprule
Model & Scales & Lang. & Dyads & Agent-games & Decisions & Cell $n$ & Coop. \\
\midrule
Claude Haiku 4.5 & 10 & 5 & 2,000 & 4,000 & 40,000 & 40 & 0.415 \\
GPT-5.4 Nano & 10 & 5 & 2,000 & 4,000 & 40,000 & 40 & 0.606 \\
Gemini 3.1 Flash-Lite & 10 & 5 & 2,000 & 4,000 & 40,000 & 40 & 0.617 \\
Gemini 3.5 Flash-Lite & 10 & 5 & 2,000 & 4,000 & 40,000 & 40 & 0.751 \\
Qwen3 235B-A22B & 10 & 5 & 2,000 & 4,000 & 40,000 & 40 & 0.435 \\
\midrule
\textbf{All} & 10 & 5 & 10,000 & 20,000 & 200,000 & 40 & 0.565 \\
\bottomrule
\end{tabular}
\end{table}

\begin{table}[htbp]
\centering
\footnotesize
\setlength{\tabcolsep}{4pt}
\caption{\textbf{Cooperation by model and payoff scale.} Each entry aggregates 200 dyads. The bracketed line under each model is a 95\% confidence interval from a bootstrap that resamples whole dyads.}
\label{tab:S-scale}
\resizebox{\textwidth}{!}{%
\begin{tabular}{@{}lcccccccccc@{}}
\toprule
Model & $\lam{=}0.01$ & $\lam{=}0.1$ & $\lam{=}0.25$ & $\lam{=}0.5$ & $\lam{=}1$ & $\lam{=}2$ & $\lam{=}5$ & $\lam{=}10$ & $\lam{=}100$ & $\lam{=}1000$ \\
\midrule
Claude Haiku 4.5 & 0.407 & 0.379 & 0.371 & 0.427 & 0.452 & 0.444 & 0.424 & 0.435 & 0.407 & 0.401 \\
 & \tiny{[0.37,0.45]} & \tiny{[0.34,0.42]} & \tiny{[0.34,0.40]} & \tiny{[0.39,0.46]} & \tiny{[0.42,0.49]} & \tiny{[0.41,0.48]} & \tiny{[0.39,0.46]} & \tiny{[0.40,0.47]} & \tiny{[0.38,0.43]} & \tiny{[0.37,0.43]} \\[2pt]
GPT-5.4 Nano & 0.438 & 0.540 & 0.644 & 0.598 & 0.576 & 0.637 & 0.680 & 0.633 & 0.663 & 0.653 \\
 & \tiny{[0.40,0.48]} & \tiny{[0.50,0.58]} & \tiny{[0.61,0.68]} & \tiny{[0.56,0.63]} & \tiny{[0.54,0.61]} & \tiny{[0.60,0.67]} & \tiny{[0.65,0.71]} & \tiny{[0.60,0.67]} & \tiny{[0.64,0.69]} & \tiny{[0.62,0.68]} \\[2pt]
Gemini 3.1 Flash-Lite & 0.641 & 0.648 & 0.621 & 0.640 & 0.635 & 0.614 & 0.614 & 0.598 & 0.584 & 0.577 \\
 & \tiny{[0.61,0.67]} & \tiny{[0.62,0.68]} & \tiny{[0.59,0.65]} & \tiny{[0.61,0.67]} & \tiny{[0.61,0.67]} & \tiny{[0.58,0.65]} & \tiny{[0.58,0.65]} & \tiny{[0.56,0.63]} & \tiny{[0.54,0.62]} & \tiny{[0.54,0.61]} \\[2pt]
Gemini 3.5 Flash-Lite & 0.792 & 0.694 & 0.872 & 0.800 & 0.774 & 0.710 & 0.697 & 0.754 & 0.734 & 0.682 \\
 & \tiny{[0.74,0.84]} & \tiny{[0.64,0.74]} & \tiny{[0.85,0.89]} & \tiny{[0.77,0.83]} & \tiny{[0.75,0.80]} & \tiny{[0.67,0.74]} & \tiny{[0.66,0.73]} & \tiny{[0.72,0.78]} & \tiny{[0.70,0.77]} & \tiny{[0.64,0.72]} \\[2pt]
Qwen3 235B-A22B & 0.476 & 0.494 & 0.455 & 0.432 & 0.416 & 0.427 & 0.419 & 0.422 & 0.399 & 0.411 \\
 & \tiny{[0.43,0.52]} & \tiny{[0.44,0.54]} & \tiny{[0.41,0.50]} & \tiny{[0.39,0.47]} & \tiny{[0.38,0.45]} & \tiny{[0.39,0.46]} & \tiny{[0.38,0.45]} & \tiny{[0.39,0.45]} & \tiny{[0.37,0.43]} & \tiny{[0.38,0.44]} \\[2pt]
\bottomrule
\end{tabular}
}
\end{table}

\begin{table}[htbp]
\centering
\scriptsize
\setlength{\tabcolsep}{2.5pt}
\caption{\textbf{Cooperation by model, prompt language and payoff scale.} 40 dyads per entry. The final column is the range across the ten scales for that model and language; it runs from 0.080 to 0.449, so the scale sensitivity of a model is a property of the model and the language jointly.}
\label{tab:S-language}
\resizebox{\textwidth}{!}{%
\begin{tabular}{@{}llccccccccccc@{}}
\toprule
Model & Language & 0.01 & 0.1 & 0.25 & 0.5 & 1 & 2 & 5 & 10 & 100 & 1000 & Range \\
\midrule
Claude Haiku 4.5 & English & 0.35 & 0.40 & 0.32 & 0.42 & 0.44 & 0.45 & 0.41 & 0.42 & 0.45 & 0.46 & \textbf{0.136} \\
 & French & 0.65 & 0.47 & 0.50 & 0.33 & 0.37 & 0.43 & 0.41 & 0.40 & 0.38 & 0.40 & \textbf{0.314} \\
 & Vietnamese & 0.17 & 0.17 & 0.28 & 0.29 & 0.26 & 0.20 & 0.30 & 0.22 & 0.23 & 0.23 & \textbf{0.130} \\
 & Chinese & 0.53 & 0.52 & 0.45 & 0.61 & 0.69 & 0.57 & 0.56 & 0.66 & 0.58 & 0.56 & \textbf{0.245} \\
 & Arabic & 0.35 & 0.33 & 0.31 & 0.48 & 0.51 & 0.56 & 0.43 & 0.47 & 0.40 & 0.36 & \textbf{0.249} \\
\addlinespace
GPT-5.4 Nano & English & 0.45 & 0.64 & 0.85 & 0.88 & 0.86 & 0.90 & 0.88 & 0.87 & 0.89 & 0.76 & \textbf{0.449} \\
 & French & 0.75 & 0.78 & 0.78 & 0.57 & 0.55 & 0.60 & 0.62 & 0.67 & 0.55 & 0.74 & \textbf{0.226} \\
 & Vietnamese & 0.49 & 0.46 & 0.58 & 0.47 & 0.50 & 0.62 & 0.75 & 0.54 & 0.65 & 0.63 & \textbf{0.289} \\
 & Chinese & 0.24 & 0.45 & 0.65 & 0.60 & 0.49 & 0.57 & 0.61 & 0.58 & 0.62 & 0.65 & \textbf{0.414} \\
 & Arabic & 0.25 & 0.37 & 0.37 & 0.48 & 0.47 & 0.49 & 0.54 & 0.51 & 0.61 & 0.49 & \textbf{0.360} \\
\addlinespace
Gemini 3.1 Flash-Lite & English & 0.74 & 0.62 & 0.67 & 0.62 & 0.63 & 0.62 & 0.59 & 0.62 & 0.57 & 0.62 & \textbf{0.167} \\
 & French & 0.57 & 0.66 & 0.70 & 0.63 & 0.65 & 0.60 & 0.60 & 0.61 & 0.60 & 0.56 & \textbf{0.145} \\
 & Vietnamese & 0.58 & 0.56 & 0.54 & 0.64 & 0.57 & 0.55 & 0.57 & 0.49 & 0.47 & 0.47 & \textbf{0.170} \\
 & Chinese & 0.62 & 0.66 & 0.53 & 0.61 & 0.59 & 0.60 & 0.60 & 0.58 & 0.58 & 0.58 & \textbf{0.133} \\
 & Arabic & 0.69 & 0.74 & 0.66 & 0.70 & 0.74 & 0.69 & 0.71 & 0.69 & 0.69 & 0.66 & \textbf{0.080} \\
\addlinespace
Gemini 3.5 Flash-Lite & English & 0.94 & 0.66 & 0.92 & 0.84 & 0.74 & 0.59 & 0.53 & 0.65 & 0.64 & 0.60 & \textbf{0.412} \\
 & French & 0.91 & 0.63 & 0.86 & 0.72 & 0.78 & 0.63 & 0.75 & 0.78 & 0.59 & 0.61 & \textbf{0.317} \\
 & Vietnamese & 0.72 & 0.63 & 0.93 & 0.86 & 0.81 & 0.76 & 0.62 & 0.64 & 0.89 & 0.73 & \textbf{0.314} \\
 & Chinese & 0.69 & 0.80 & 0.76 & 0.64 & 0.70 & 0.73 & 0.67 & 0.76 & 0.66 & 0.73 & \textbf{0.162} \\
 & Arabic & 0.70 & 0.75 & 0.89 & 0.94 & 0.83 & 0.84 & 0.92 & 0.93 & 0.89 & 0.74 & \textbf{0.242} \\
\addlinespace
Qwen3 235B-A22B & English & 0.47 & 0.51 & 0.53 & 0.44 & 0.42 & 0.49 & 0.43 & 0.50 & 0.37 & 0.46 & \textbf{0.155} \\
 & French & 0.41 & 0.40 & 0.35 & 0.39 & 0.37 & 0.36 & 0.35 & 0.35 & 0.34 & 0.33 & \textbf{0.083} \\
 & Vietnamese & 0.52 & 0.55 & 0.48 & 0.48 & 0.45 & 0.45 & 0.47 & 0.44 & 0.40 & 0.42 & \textbf{0.150} \\
 & Chinese & 0.48 & 0.52 & 0.46 & 0.44 & 0.43 & 0.45 & 0.45 & 0.44 & 0.49 & 0.47 & \textbf{0.084} \\
 & Arabic & 0.50 & 0.50 & 0.46 & 0.42 & 0.40 & 0.40 & 0.40 & 0.38 & 0.39 & 0.38 & \textbf{0.124} \\
\addlinespace
\bottomrule
\end{tabular}
}
\end{table}

\begin{table}[htbp]
\centering
\footnotesize
\setlength{\tabcolsep}{4pt}
\caption{\textbf{Variance decomposition of the cell means.} Type-II analysis of variance on the 250 model $\times$ language $\times$ scale cell means, where the design is exactly balanced. The payoff scale has no significant main effect while both of its interactions are significant.}
\label{tab:S-variance}
\begin{tabular}{@{}lrrrrr@{}}
\toprule
Term & Sum sq. & df & $F$ & $p$ & \% variance \\
\midrule
Model & 3.9188 & 4 & 206.80 & $<0.001$ & 56.5 \\
Language & 0.2092 & 4 & 11.04 & $<0.001$ & 3.0 \\
Payoff scale & 0.0471 & 9 & 1.10 & $0.364$ & 0.7 \\
Model $\times$ language & 1.3665 & 16 & 18.03 & $<0.001$ & 19.7 \\
Model $\times$ scale & 0.4542 & 36 & 2.66 & $<0.001$ & 6.5 \\
Language $\times$ scale & 0.2615 & 36 & 1.53 & $0.041$ & 3.8 \\
Residual & 0.6822 & 144 &  &  & 9.8 \\
\bottomrule
\end{tabular}
\end{table}

\begin{table}[htbp]
\centering
\footnotesize
\setlength{\tabcolsep}{4pt}
\caption{\textbf{Persona effect by model and payoff scale.} The cooperation rate of an agent instructed that it is cooperative minus that of an agent instructed that it is selfish, at each scale. The two leftmost columns are the scales at which every payoff printed is at most 1.}
\label{tab:S-persona}
\resizebox{\textwidth}{!}{%
\begin{tabular}{@{}lcccccccccc@{}}
\toprule
Model & $\lam{=}0.01$ & $\lam{=}0.1$ & $\lam{=}0.25$ & $\lam{=}0.5$ & $\lam{=}1$ & $\lam{=}2$ & $\lam{=}5$ & $\lam{=}10$ & $\lam{=}100$ & $\lam{=}1000$ \\
\midrule
Claude Haiku 4.5 & $-0.211$ & $-0.232$ & $+0.206$ & $+0.200$ & $+0.179$ & $+0.261$ & $+0.277$ & $+0.158$ & $+0.168$ & $+0.211$ \\
GPT-5.4 Nano & $-0.172$ & $-0.098$ & $+0.015$ & $-0.026$ & $-0.053$ & $+0.032$ & $+0.041$ & $+0.055$ & $-0.017$ & $+0.100$ \\
Gemini 3.1 Flash-Lite & $+0.552$ & $+0.602$ & $+0.640$ & $+0.670$ & $+0.693$ & $+0.681$ & $+0.707$ & $+0.696$ & $+0.659$ & $+0.692$ \\
Gemini 3.5 Flash-Lite & $-0.335$ & $-0.578$ & $+0.122$ & $+0.234$ & $+0.241$ & $+0.258$ & $+0.313$ & $+0.203$ & $+0.193$ & $+0.230$ \\
Qwen3 235B-A22B & $-0.901$ & $-0.922$ & $-0.811$ & $-0.768$ & $-0.702$ & $-0.635$ & $-0.690$ & $-0.602$ & $-0.657$ & $-0.566$ \\
\bottomrule
\end{tabular}
}
\end{table}

\begin{table}[htbp]
\centering
\footnotesize
\setlength{\tabcolsep}{4pt}
\caption{\textbf{Share of agent-games matched exactly by one canonical rule (\%).} Computed by exact rule matching, with no learned component. 400 agent-games per entry.}
\label{tab:S-deduced}
\resizebox{\textwidth}{!}{%
\begin{tabular}{@{}lcccccccccc@{}}
\toprule
Model & $\lam{=}0.01$ & $\lam{=}0.1$ & $\lam{=}0.25$ & $\lam{=}0.5$ & $\lam{=}1$ & $\lam{=}2$ & $\lam{=}5$ & $\lam{=}10$ & $\lam{=}100$ & $\lam{=}1000$ \\
\midrule
Claude Haiku 4.5 & 20.0 & 15.8 & 3.2 & 3.5 & 4.0 & 2.8 & 4.0 & 2.5 & 4.5 & 3.5 \\
GPT-5.4 Nano & 12.5 & 7.0 & 9.2 & 6.2 & 5.5 & 6.5 & 8.2 & 7.2 & 8.5 & 9.0 \\
Gemini 3.1 Flash-Lite & 16.5 & 19.0 & 21.2 & 23.5 & 29.0 & 23.0 & 30.0 & 29.5 & 27.5 & 26.8 \\
Gemini 3.5 Flash-Lite & 30.8 & 33.8 & 20.2 & 13.0 & 20.5 & 14.2 & 14.8 & 11.5 & 14.0 & 16.2 \\
Qwen3 235B-A22B & 77.0 & 67.2 & 63.0 & 62.7 & 59.8 & 45.2 & 58.0 & 43.2 & 55.0 & 42.0 \\
\bottomrule
\end{tabular}
}
\end{table}

\begin{table}[htbp]
\centering
\scriptsize
\setlength{\tabcolsep}{2.5pt}
\caption{\textbf{Strategy composition by model and payoff scale (\%).} The label assigned by the hybrid read-out: the rule base fixes the candidate set and the LSTM ranks within it, deciding alone only where no canonical rule fits. Columns sum to 100 within each model and scale.}
\label{tab:S-mix}
\resizebox{\textwidth}{!}{%
\begin{tabular}{@{}llcccccccccc@{}}
\toprule
Model & Label & 0.01 & 0.1 & 0.25 & 0.5 & 1 & 2 & 5 & 10 & 100 & 1000 \\
\midrule
Claude Haiku 4.5 & AllC & 27.0 & 20.0 & 14.2 & 20.0 & 21.8 & 18.2 & 15.2 & 18.8 & 12.8 & 14.0 \\
 & TFT & 20.5 & 23.5 & 23.5 & 18.0 & 16.8 & 16.5 & 19.5 & 17.8 & 20.8 & 22.2 \\
 & WSLS & 8.8 & 13.5 & 15.2 & 13.5 & 20.2 & 21.2 & 20.5 & 18.2 & 21.8 & 19.0 \\
 & AllD & 43.8 & 43.0 & 47.0 & 48.5 & 41.2 & 44.0 & 44.8 & 45.2 & 44.8 & 44.8 \\
\addlinespace
GPT-5.4 Nano & AllC & 29.8 & 41.8 & 52.2 & 46.2 & 48.5 & 54.2 & 58.5 & 52.2 & 53.2 & 50.7 \\
 & TFT & 11.8 & 11.5 & 10.5 & 11.8 & 10.2 & 13.2 & 13.8 & 13.2 & 13.0 & 17.8 \\
 & WSLS & 12.5 & 12.5 & 14.5 & 15.5 & 13.8 & 10.5 & 12.5 & 10.2 & 15.5 & 14.0 \\
 & AllD & 46.0 & 34.2 & 22.8 & 26.5 & 27.5 & 22.0 & 15.2 & 24.2 & 18.2 & 17.5 \\
\addlinespace
Gemini 3.1 Flash-Lite & AllC & 55.8 & 57.0 & 58.2 & 61.0 & 59.8 & 55.0 & 58.5 & 53.8 & 52.2 & 52.2 \\
 & TFT & 6.5 & 3.8 & 2.2 & 0.2 & 2.5 & 3.5 & 1.2 & 2.2 & 6.2 & 5.8 \\
 & WSLS & 12.8 & 13.2 & 8.8 & 2.8 & 8.0 & 7.2 & 6.5 & 7.2 & 7.2 & 6.0 \\
 & AllD & 25.0 & 26.0 & 30.8 & 36.0 & 29.8 & 34.2 & 33.8 & 36.8 & 34.2 & 36.0 \\
\addlinespace
Gemini 3.5 Flash-Lite & AllC & 73.8 & 59.8 & 74.0 & 64.0 & 58.0 & 60.0 & 51.7 & 59.8 & 56.2 & 51.5 \\
 & TFT & 6.8 & 2.2 & 10.2 & 16.2 & 17.0 & 16.2 & 15.2 & 13.5 & 15.2 & 14.2 \\
 & WSLS & 6.2 & 13.2 & 14.0 & 16.2 & 21.2 & 12.0 & 18.0 & 18.5 & 16.5 & 15.2 \\
 & AllD & 13.2 & 24.8 & 1.8 & 3.5 & 3.8 & 11.8 & 15.0 & 8.2 & 12.0 & 19.0 \\
\addlinespace
Qwen3 235B-A22B & AllC & 42.0 & 47.8 & 39.2 & 38.0 & 36.8 & 36.5 & 35.5 & 35.2 & 34.5 & 33.8 \\
 & TFT & 0.5 & 0.0 & 0.0 & 0.0 & 1.2 & 0.8 & 2.0 & 6.2 & 0.0 & 4.8 \\
 & WSLS & 9.0 & 2.5 & 8.8 & 13.5 & 7.2 & 12.0 & 8.5 & 8.2 & 9.2 & 10.0 \\
 & AllD & 48.5 & 49.8 & 52.0 & 48.5 & 54.8 & 50.7 & 54.0 & 50.2 & 56.2 & 51.5 \\
\addlinespace
\bottomrule
\end{tabular}
}
\end{table}

\begin{table}[htbp]
\centering
\footnotesize
\setlength{\tabcolsep}{4pt}
\caption{\textbf{Provenance of each label.} For every label, the share of the agent-games carrying it that were deduced (exactly one canonical rule fits), ambiguous (several fit, and the LSTM ranked within that set) or unmatched (no canonical rule fits, so the label is an attribution the LSTM makes alone). The asymmetry is severe and bounds how the compositional result may be read: AllC and AllD carry a substantial deduced share, whereas the TFT and WSLS labels are almost entirely attributions.}
\label{tab:S-provenance}
\begin{tabular}{@{}lrccc@{}}
\toprule
Label & $n$ & Deduced (\%) & Ambiguous (\%) & Unmatched (\%) \\
\midrule
AllC & 8,884 & 26.0 & 29.6 & 44.4 \\
TFT & 2,010 & 4.6 & 0.0 & 95.4 \\
WSLS & 2,494 & 2.0 & 0.7 & 97.3 \\
AllD & 6,612 & 32.7 & 0.0 & 67.3 \\
\bottomrule
\end{tabular}
\end{table}

\begin{table}[htbp]
\centering
\footnotesize
\setlength{\tabcolsep}{4pt}
\caption{\textbf{Validation of the LSTM branch.} Trained on synthetic AllC / AllD / Tit-for-Tat / Win-Stay-Lose-Shift trajectories at zero and 5\% execution noise, balanced at 40{,}320 sequences per strategy per level. The 10\% and 20\% noise levels are never trained on. No language-model transcript enters training.}
\label{tab:S-classifier}
\begin{tabular}{@{}lrc@{}}
\toprule
Split & $n$ & Accuracy \\
\midrule
held-out (NoNoise + Noise005) & 64,512 & 0.9846 \\
OOD (Noise01) & 161,280 & 0.9364 \\
OOD (Noise02) & 161,280 & 0.8107 \\
held-out recall: AllC & 16,150 & 0.9853 \\
held-out recall: AllD & 16,203 & 0.9952 \\
held-out recall: TFT & 16,193 & 0.9744 \\
held-out recall: WSLS & 15,966 & 0.9834 \\
\bottomrule
\end{tabular}
\end{table}

\begin{table}[htbp]
\centering
\footnotesize
\setlength{\tabcolsep}{4pt}
\caption{\textbf{Pairwise correlation between the model curves.} Correlation of the ten-point cooperation-against-scale curves of each pair of models. Six of the ten pairs are negative and the mean is $-0.10$, which is why averaging the curves cancels most of the movement.}
\label{tab:S-shapes}
\begin{tabular}{@{}lccccc@{}}
\toprule
 & Claude Haiku & GPT-5.4 Nano & Gemini 3.1 & Gemini 3.5 & Qwen3 235B-A22B \\
\midrule
Claude Haiku 4.5 &  & $+0.11$ & $-0.03$ & $-0.18$ & $-0.56$ \\
GPT-5.4 Nano & $+0.11$ &  & $-0.70$ & $-0.23$ & $-0.72$ \\
Gemini 3.1 Flash-Lite & $-0.03$ & $-0.70$ &  & $+0.36$ & $+0.73$ \\
Gemini 3.5 Flash-Lite & $-0.18$ & $-0.23$ & $+0.36$ &  & $+0.22$ \\
Qwen3 235B-A22B & $-0.56$ & $-0.72$ & $+0.73$ & $+0.22$ &  \\
\bottomrule
\end{tabular}
\end{table}


\end{document}
